\section{Methods}
\label{sec:methods}

We organise models as a ladder (Stages~0--5; Figure~\ref{fig:pipeline}). Two
invariants make the comparison honest. First, from Stage~2 upward every learned
model includes past-volatility covariates ($v_{\mathrm{pre}}$ at all horizons plus
the HAR-RV forecast) and is \emph{also} reported with them ablated, baking the
persistence control into every comparison. Second, given the small-data regime
($10^3$--$10^4$ samples) we use frozen representations with light heads (ridge,
shallow MLP, LightGBM) and per-horizon models by default; fine-tuning is gated to
Stage~6 and a shared multi-task head is an ablation, not the headline. Because the
shallow-MLP heads proved seed-unstable (across-seed \Rtwo{} std up to $\approx$3.4;
see Limitations), the consolidated grid (Table~\ref{tab:grid}) uses one
pre-registered deterministic ridge/structural canonical model per stage; MLP heads
appear in the appendix only.

% Stage-ladder schematic (Stage 0 -> 4, plus the reproduction track). Simple TikZ
% stub; refine for camera-ready.
\begin{figure}[t]
\centering
\begin{tikzpicture}[
    node distance=3.2mm,
    box/.style={draw, rounded corners, align=center, font=\scriptsize,
                inner sep=2pt, text width=0.9\linewidth, minimum height=6mm},
    done/.style={box, fill=green!8},
    todo/.style={box, fill=orange!10},
]
\node[done] (s0) {\textbf{Stage 0} Persistence, EWMA, HAR-RV, GARCH(1,1)};
\node[done, below=of s0] (s1) {\textbf{Stage 1} LightGBM $+$ ticker fixed effect};
\node[done, below=of s1] (s2) {\textbf{Stage 2} Frozen text: BGE-M3, FinBERT, surface $\to$ ridge / MLP};
\node[done, below=of s2] (s3) {\textbf{Stage 3} Frozen audio: eGeMAPS, WavLM, emotion2vec+};
\node[done, below=of s3] (s4) {\textbf{Stage 4} Fusion (gated / stacking)};
\node[todo, below=of s4] (sr) {\textbf{Stage R} Prior models (HTML, SCSS, KeFVP, DialogueGAT, VolTAGE) on their own data, under the same controls};
\foreach \a/\b in {s0/s1, s1/s2, s2/s3, s3/s4, s4/sr}
  \draw[->, >=stealth] (\a) -- (\b);
\end{tikzpicture}
\caption{\textbf{The modelling ladder.} From Stage~2 up, every model also carries
past-volatility covariates (and is reported with them ablated). Every stage runs
through the identity-control suite. Green = complete, orange = pending.}
\label{fig:pipeline}
\end{figure}


\paragraph{Stage 0 --- econometric baselines.}
Persistence ($\hat v = v_{\mathrm{pre}}$), EWMA (RiskMetrics $\lambda{=}0.94$),
HAR-RV (OLS on daily/weekly/monthly realised-variance factors,
\citealp{corsi2009har}), and GARCH(1,1) (per-call fit, \texttt{arch}). These
establish the floor and the denominator of the headline metric.

\paragraph{Stage 1 --- tabular identity baseline.}
LightGBM over past-volatility features, trivial metadata (transcript length, Q\&A
turn count where available), and a \textbf{ticker fixed effect}. This is the
in-house ``Same-Company-Same-Signal'' control: it measures how much apparent
multimodal performance is identity plus metadata.

\paragraph{Stage 2 --- frozen text content.}
We segment each transcript into prepared-remarks vs.\ Q\&A and speaker turns, then
extract three feature blocks: BGE-M3 \citep{bge2024} section-pooled sentence
embeddings (1024-d), FinBERT \citep{finbert} sentiment aggregates per
scope$\times$role, and surface statistics. Heads are ridge (with the $\alpha$ grid
selected on validation, \emph{no} retrain on train$+$val) and a shallow MLP
(1$\times$256, L2 $+$ early stopping, 5 seeds). Features use fp32 with content-hash
caching for bit-identical re-extraction.

\paragraph{Stage 3 --- frozen audio content.}
Audio is QC'd and resampled to 16\,kHz mono FLAC. We extract eGeMAPSv02 functionals
(88-d, openSMILE \citep{opensmile}, CPU, interpretable), then WavLM-Large
\citep{wavlm} and emotion2vec+ \citep{emotion2vec} embeddings over 30\,s
non-overlapping windows, mean-pooled to a per-call vector (2{,}671 calls, all
extracted on one consumer GPU with resumable content-hash caches), fed to the same
heads. We use per-call pooling only (no per-sentence alignment), and report a
gender-confound analysis.

\paragraph{Stage 4 --- fusion.}
Gated fusion (per-modality train-fit PCA + norm gates $\to$ shallow MLP) and
late-fusion stacking (a meta-ridge over the text, audio, and Stage-1 GBDT base
predictions, fit on validation only) over frozen modality embeddings, to test
whether multimodal beats the best unimodal model (RQ2). Cross-attention is omitted:
with one pooled vector per modality, attention over 2--3 tokens reduces to gated
fusion. All fusion heads are $<$5M parameters.

\paragraph{Stage 5 --- LLM-extracted structured features.}
Qwen2.5-7B-Instruct \citep{qwen25} (4-bit) with grammar-constrained JSON decoding
(JSON-schema grammar via \texttt{llama.cpp}; Outlines \citep{outlines} on the
GPU-server path) extracts auditable per-section features (guidance direction,
hedging intensity, Q\&A evasiveness, surprise mentions, analyst tone) against a
frozen, human-signed-off schema and rubric, with YaRN-extended 65{,}536-token
context so whole sections are never truncated. A pre-registered human-audit gate
governs the stage: corpus-scale extraction is blocked until the model reaches
$\kappa>0.6$ against human rubric labels on 50 frozen train-only calls, extracted
through the \emph{same} engine and configuration as the corpus run (RQ3). As
\S\ref{sec:results-llm} reports, the gate failed and corpus extraction was not
performed.
